% !TeX program = tectonic
\documentclass[11pt,a4paper]{article}
\usepackage{fontspec}
\usepackage[english]{babel}
\babelfont{rm}[Language=Default]{Liberation Serif}
\babelfont[Language=Default]{sf}{Carlito}
\babelfont[Language=Default]{tt}{DejaVu Sans Mono}
\usepackage{amsmath,amssymb,amsthm}
\usepackage{geometry}
\geometry{a4paper, top=2.4cm, bottom=2.4cm, left=2.4cm, right=2.4cm}
\usepackage{booktabs,tabularx,enumitem,xcolor,tcolorbox}
\tcbuselibrary{breakable,skins}
\definecolor{linkblue}{HTML}{1F4E79}
\definecolor{statgreen}{HTML}{2F6B3A}
\definecolor{goldacc}{HTML}{8B7E5A}
\usepackage[colorlinks=true,linkcolor=linkblue,urlcolor=linkblue,unicode=true]{hyperref}
\theoremstyle{plain}
\newtheorem{theorem}{Theorem}
\newtheorem{lemma}[theorem]{Lemma}
\newtheorem{proposition}[theorem]{Proposition}
\newtheorem{corollary}[theorem]{Corollary}
\theoremstyle{definition}
\newtheorem{definition}[theorem]{Definition}
\theoremstyle{remark}
\newtheorem{remark}[theorem]{Remark}
\newtcolorbox{statusbox}[1][]{colback=statgreen!5,colframe=statgreen!75,
  title={\textbf{Summary of results:} #1},fonttitle=\small,boxrule=0.7pt,arc=2pt,breakable}
\newtcolorbox{databox}[1][]{colback=goldacc!6,colframe=goldacc!80,
  title={\textbf{Protocol data:} #1},fonttitle=\small,boxrule=0.7pt,arc=2pt,breakable}
\newtcolorbox{verifybox}[1][]{colback=linkblue!5,colframe=linkblue!80,
  title={\textbf{Verification:} #1},fonttitle=\small,boxrule=0.7pt,arc=2pt,breakable}
\widowpenalty=10000
\clubpenalty=10000
\setlength{\parskip}{0.35em}
\newcommand{\CC}{\mathbb{C}}
\newcommand{\RR}{\mathbb{R}}
\newcommand{\ZZ}{\mathbb{Z}}
\newcommand{\QQ}{\mathbb{Q}}
\newcommand{\Ga}{\Gamma}
\newcommand{\Bfun}{\mathrm{B}}
\newcommand{\im}{\operatorname{Im}}
\newcommand{\re}{\operatorname{Re}}
\newcommand{\lcm}{\operatorname{lcm}}
\newcommand{\SNF}{\operatorname{SNF}}
\newcommand{\Jac}{\operatorname{Jac}}
\newcommand{\rank}{\operatorname{rank}}
\newcommand{\Ind}{\operatorname{Index}}
\newcommand{\disc}{\operatorname{disc}}
\begin{document}
\begin{center}
{\LARGE\bfseries Certificate C: Mu4-Equivariance of the Corrections}\\[0.4em]
{\large 22 = 1 + 7 + 7 + 7 and the regular representation}\\[0.6em]
{\normalsize Isaev Iskhak Khamzatovich}\\[0.2em]
{\normalsize The <<Dynamic Principle>> Program --- the \texttt{hodge-laboratory}}\\[0.15em]
{\small Standalone theorem monograph · Монография 11/16}
\end{center}
\vspace{0.4em}
{\color{linkblue}\hrule height 0.8pt}
\vspace{0.8em}

\section{Setting}

Certificate C establishes the naturality of the corrections: the
triple $(7,22,\lambda_1)$ is consistent with the $\mu_4$
representation --- the group of quarter-turns present on all three
stands.

\begin{theorem}[C1 --- K3 isotypy]\label{th:C1stand}
In $H^2$ of the Fermat quartic the isotypic decomposition over the
$\mu_4$ characters is $22=1+7+7+7$. Three independent counting
routes (direct orbits, the spectrum of the permutation
representation, exact character multiplicities) agree.
\end{theorem}

\begin{proof}
The action of $(\mu_4)^2$ on the $48$ lines has $12$ orbits. The
permutation representation decomposes over the characters;
projection operators $\Pi_\chi=\frac14\sum_k i^{-k}R^k$ extract the
isotypic parts. The orbit count gives $(1,7,7,7)$; the spectrum of
the permutation representation gives the same multiplicities as
eigenspace dimensions; the pairing with exact line characters is the
third route. Three agreeing routes exclude coincidence.
\end{proof}

\begin{theorem}[C2 --- torus, regular representation]\label{th:C2stand}
On the torus the correction operator $\lambda_1$ commutes with
$\mu_4$; the characteristic polynomial is $x^4-1$ --- the regular
representation; the multiplicity of $\lambda_1$ is $4$; the
continuum is $\lambda_1=4\pi^2$; the triple $(4,1,4\pi^2)$ is
consistent.
\end{theorem}

\begin{proof}
Commutativity --- the equivariance of the Laplacian; the eigenspaces
are invariant and decompose over the $\mu_4$ characters. The
polynomial $x^4-1$ has simple roots $\{1,-1,i,-i\}$ --- each
character exactly once: the regular representation. The value
$4\pi^2$ is realized on every character --- the lattice continuum.
\end{proof}

\begin{databox}[the Klein heptad]
$\Phi_7$ --- the cyclotomic polynomial; $\Delta_{49a1}=-343=-7^3$;
$j=-3375=-15^3$; the CM order $\ZZ[(1+\sqrt{-7})/2]$; $h(-7)=1$;
$\tau$ certified. The heptad $\sqrt{-7}$ is the exact integer
environment of the phase $\pi/7$. The $\mu_4$ stand: anomaly $=i$;
$42$ elements of order $4$ in the dimension-32 permutation
representation; $\mu_4$ on all three stands. Grid: symmetric
violations $0$; $(7,22)$ violations $1023/1024$.
\end{databox}

\begin{statusbox}[summary]
Certificate C accepted ($21/21$ checks): the isotypy
$22=1+7+7+7$ confirmed by three routes; the torus gives the regular
representation; the seven receives its exact arithmetic face.
\end{statusbox}

\begin{verifybox}[reproduction]
\texttt{python3 laboratory.py} $\to$ ``Certificates $\to$ C'':
projection operators, orbits, spectrum --- all exact.
\end{verifybox}

\vfill
{\color{linkblue}\hrule height 0.6pt}
\vspace{0.5em}
{\small\color{gray}
License: individual exclusive license of Isaev Iskhak Khamzatovich.
All rights to the computations, formulas, and texts belong to the author.
Verification: \texttt{python3 laboratory.py} (``Stands''/``Certificates''),
Lean: \texttt{verification/lean/HodgeLaboratory.lean}.
}
\end{document}
